\documentclass[11pt,a4paper]{article}

% ============================================================
% Paper 4 (versi Bahasa Indonesia) --- draft kerja untuk promotor.
% Judul: Federated Cross-Network NIDS via Semantic Feature Mapping.
% Fondasi: Paper 1 (SFM) \cite{sfm}; sejajar Paper 2 (adversarial), Paper 3 (online).
% CATATAN KEJUJURAN: seluruh ANGKA HASIL masih [TBD] -- diisi dari eksperimen NYATA
%   (notebook federated/notebooks/nb01-nb07 + eksekusi AWS). JANGAN mengarang angka.
% ============================================================

\usepackage[utf8]{inputenc}
\usepackage[T1]{fontenc}
\usepackage[bahasa]{babel}
\usepackage{graphicx}
\usepackage{algorithm}
\usepackage{algpseudocode}
\usepackage{booktabs}
\usepackage{amsmath,amssymb}
\usepackage{siunitx}
\usepackage{hyperref}
\usepackage{geometry}
\usepackage{caption}
\usepackage{multirow}
\usepackage{placeins}
\geometry{margin=2.5cm}
\graphicspath{{figure-fed/}}
\hypersetup{colorlinks=true, linkcolor=blue, citecolor=blue, urlcolor=blue}

% Penanda placeholder hasil yang belum ada (kejujuran data).
\newcommand{\TBD}{\textit{[TBD]}}

\title{\textbf{NIDS Terfederasi Lintas-Jaringan melalui Pemetaan Fitur Semantik:
Dapatkah Ruang Fitur Semantik Bersama Mengatasi Heterogenitas Skema-Fitur dalam
\emph{Federated Learning}?}}

\author{Hero Yudo Martono\thanks{Penulis koresponden: \texttt{hero@pens.ac.id}},
\quad Iwan Syarif, \quad Ferry Astika Saputra \\[0.3em]
\textit{Departemen Teknik Informatika, Politeknik Elektronika Negeri Surabaya,
Surabaya, Indonesia} \\
\texttt{hero@pens.ac.id; iwanarif@pens.ac.id; ferryas@pens.ac.id}}

\date{Draft --- \today}

\begin{document}
\maketitle

\begin{abstract}
\emph{Federated learning} (FL) menjanjikan kolaborasi antar-organisasi untuk
melatih sistem deteksi intrusi jaringan (NIDS) tanpa memusatkan data mentah yang
sensitif. Namun hampir seluruh pekerjaan FL-NIDS mengasumsikan setiap klien
memakai ruang fitur yang identik --- asumsi yang runtuh di dunia nyata karena
tiap organisasi lazim memakai \emph{alat ekstraksi fitur} berbeda (mis.
CICFlowMeter versus Argus/Bro), sehingga \emph{skema fitur} antar-klien berbeda.
Makalah ini mengusulkan \emph{Semantic Feature Mapping} (SFM) --- ruang fitur
semantik yang tervalidasi statistik dari karya pendahulu kami \cite{sfm} ---
sebagai \emph{enabler} federasi antar-klien berskema-fitur berbeda. Kontribusi
utama bukan menambahkan FL pada \emph{pipeline} lama, melainkan menunjukkan bahwa
SFM \emph{menyelesaikan} heterogenitas skema-fitur yang selama ini menggagalkan FL
lintas-\emph{dataset}: kamus pemetaan SFM diperlakukan sebagai \emph{kesepakatan
skema publik} (bukan hasil belajar dari data gabungan), diterapkan tiap klien
secara lokal dan mandiri, sehingga hanya bobot model yang dipertukarkan. Kami
mengevaluasi pada CSE-CIC-IDS2018 dan UNSW-NB15 dengan protokol multi-\emph{seed}
(MCC sebagai metrik utama, dilengkapi \emph{recall}/\emph{precision}/F1 dan
\emph{balanced accuracy}), membandingkan model terpusat (\emph{upper bound}),
lokal-saja (\emph{lower bound}), dan federated-SFM (FedAvg serta FedProx), beserta
ablasi FL dengan dan tanpa SFM dan biaya komunikasi. \textbf{Catatan:} seluruh
angka hasil pada draf ini masih \TBD{} karena eksperimen belum dieksekusi; nilai
akan diisi dari hasil nyata sebelum diajukan untuk publikasi.

\medskip
\noindent\textbf{Kata kunci:} deteksi intrusi jaringan; \emph{federated learning};
pemetaan fitur semantik; heterogenitas skema-fitur; interoperabilitas; privasi.
\end{abstract}

% ============================================================
\section{Pendahuluan}
\label{sec:intro}

Sistem deteksi intrusi jaringan berbasis pembelajaran mesin mencapai akurasi
tinggi pada pengaturan \emph{dalam-jaringan}, tetapi kolaborasi antar-organisasi
untuk berbagi pengetahuan deteksi terbentur masalah privasi dan regulasi: data
trafik mentah jarang boleh keluar dari batas organisasi. \emph{Federated
learning} \cite{popoola} menawarkan jalan keluar --- melatih model bersama hanya
dengan bertukar parameter, bukan data. Dalam ranah NIDS, FL telah banyak dikaji
dari sisi arsitektur, agregasi, non-IID, hingga ketahanan \emph{Byzantine}
\cite{fsurvey1,fsurvey2,ftaxonomy}.

Namun terdapat celah mendasar yang jarang disentuh. Mayoritas kerangka FL-NIDS
mengandaikan seluruh klien berbagi \emph{ruang fitur yang sama}. Di lapangan,
asumsi ini tidak realistis: organisasi berbeda memakai perangkat ekstraksi fitur
berbeda, dengan penamaan dan definisi kolom yang tidak sepadan. Satu klien
mungkin menghasilkan fitur ala CICFlowMeter, klien lain ala Argus atau Zeek.
Akibatnya, FedAvg \cite{popoola} --- yang merata-ratakan bobot per-dimensi ---
menjadi tak-bermakna karena dimensi ke-$i$ pada satu klien tidak sepadan secara
semantik dengan dimensi ke-$i$ pada klien lain. Kami menyebut persoalan ini
\emph{heterogenitas skema-fitur}.

\paragraph{Pertanyaan riset.} Dapatkah \emph{Semantic Feature Mapping} (SFM) ---
ruang fitur semantik yang tervalidasi statistik dari karya pendahulu kami
\cite{sfm} --- berperan sebagai \emph{enabler} federasi antar-klien
berskema-fitur berbeda, dan seberapa jauh federated-SFM menutup celah performa
terhadap model terpusat tanpa memusatkan data mentah?

\paragraph{Kontribusi.} Makalah ini menyumbang:
\begin{enumerate}
  \item \textbf{SFM sebagai jembatan interoperabilitas fitur dalam FL.} Untuk
  pertama kalinya, ruang fitur semantik tervalidasi statistik dipakai agar klien
  berskema berbeda (CIC versus UNSW) dapat berfederasi dalam satu model global
  tanpa berbagi data mentah. Kebaruan terletak pada peran SFM sebagai
  \emph{penyelesai} heterogenitas skema, bukan sekadar penambahan modul FL.
  \item \textbf{Kerangka privasi yang eksplisit.} Kamus SFM diperlakukan sebagai
  \emph{public prior} (kesepakatan skema), diterapkan tiap klien secara lokal;
  normalisasi dilakukan per-klien. Kami rinci apa yang \emph{tidak} dilakukan
  (tak ada \emph{fitting} pada data gabungan) untuk menutup celah kebocoran.
  \item \textbf{Evaluasi menyeluruh dan jujur.} Perbandingan terpusat / lokal-saja
  / federated-SFM, ablasi dengan--tanpa SFM, pengaruh non-IID dan jumlah klien,
  FedAvg versus FedProx, serta biaya komunikasi, dengan protokol multi-\emph{seed}
  dan uji signifikansi (pola Paper~2).
\end{enumerate}

\paragraph{Posisi terhadap karya kami sebelumnya (anti \emph{salami-slicing}).}
Paper~1 \cite{sfm} memperkenalkan SFM + \emph{few-shot} untuk generalisasi
lintas-jaringan pada model \emph{terpusat}. Paper~2 membahas ketahanan
adversarial, juga terpusat. Paper~3 menangani adaptasi \emph{online} terhadap
pergeseran distribusi. Makalah ini (Paper~4) membuka sumbu baru:
\emph{desentralisasi/privasi}. SFM dipindahkan ke pengaturan \emph{federated};
kontribusi orisinalnya adalah SFM sebagai \emph{enabler} federasi lintas-skema,
diukur terhadap batas terpusat dan lokal-saja. Dimensi \emph{few-shot} dan
adversarial sengaja tidak disertakan agar klaim tetap tajam.

% ============================================================
\FloatBarrier
\section{Karya Terkait}
\label{sec:related}

\paragraph{FL untuk NIDS.} Survei dan taksonomi terkini
\cite{fsurvey1,fsurvey2,ftaxonomy,falgo} memetakan arsitektur, strategi agregasi,
dan tantangan terbuka FL-NIDS. Sejumlah kerangka menekankan efisiensi sumber daya
\cite{fresource,fdbscan}, pengaturan \emph{cross-silo} berbasis graf \cite{fentente},
atau IIoT berskala besar \cite{fiiot}. Karya-karya ini memperkuat kelayakan FL,
tetapi hampir seluruhnya mengandaikan ruang fitur homogen antar-klien.

\paragraph{Non-IID dan heterogenitas.} Heterogenitas \emph{statistik} (distribusi
data berbeda antar-klien) telah banyak dikaji \cite{fnoniid,fpersonal,fmist}, dan
batas FL untuk IDS mulai disorot secara kritis \cite{flimit}. Namun heterogenitas
yang kami angkat berbeda jenis: bukan sekadar distribusi berbeda pada ruang fitur
sama, melainkan \emph{ruang fitur yang berbeda} akibat alat ekstraksi berbeda.
Upaya penyelarasan lintas-domain \cite{fcrossvan} umumnya bekerja pada ruang yang
sudah sepadan; kami justru menyediakan sepadan-nya itu melalui SFM.

\paragraph{Robustness, few-shot, dan transfer.} Garis riset ketahanan adversarial
dan \emph{Byzantine}/\emph{poisoning} dalam FL \cite{frobust,fbackdoor,fdelve}
serta \emph{few-shot}/\emph{meta}/\emph{transfer} \cite{ffewshot,ffecograph,ftransfer}
bersifat komplementer terhadap kerja kami dan dicadangkan sebagai arah lanjutan.
Fokus Paper~4 adalah \emph{interoperabilitas fitur} sebagai prasyarat yang selama
ini terlewat.

\paragraph{Set fitur standar.} Upaya menuju set fitur NIDS standar
\cite{sarhan,sarhannf} sejalan semangat kami, namun SFM menekankan \emph{validasi
semantik dan statistik} pemetaan antar-\emph{tool} serta penerapannya sebagai
\emph{public prior} dalam federasi --- bukan sekadar menyepakati nama kolom.

% ============================================================
\FloatBarrier
\section{Metode}
\label{sec:method}

\subsection{Semantic Feature Mapping sebagai \emph{public prior}}
\label{sec:sfm}

SFM adalah \emph{kamus pemetaan} yang menerjemahkan keluaran beragam alat
ekstraksi fitur ke ruang bersama berdimensi sembilan. Kesembilan fitur ---
\emph{duration}, \emph{fwd\_pkts}, \emph{bwd\_pkts}, \emph{fwd\_bytes},
\emph{bwd\_bytes}, \emph{fwd\_mean}, \emph{bwd\_mean}, \emph{src\_load},
\emph{dst\_load} --- diwarisi dari Paper~1 \cite{sfm}, di mana pemetaannya
ditemukan sekali secara \emph{offline} dengan membandingkan \emph{definisi/semantik
kolom} tiap alat (mis.\ ``Flow Duration'' pada CICFlowMeter $\approx$ ``dur'' pada
Argus) dan divalidasi secara statistik \emph{pada tiap dataset secara terpisah}.
Hasilnya berupa metadata beberapa baris, yakni pengetahuan domain --- bukan data.

Poin yang paling mungkin dipertanyakan penilai adalah: dari mana sembilan fitur
itu berasal tanpa klien berbagi data? Jawaban kami tegas: \textbf{kamus SFM adalah
kesepakatan skema (public prior), bukan hasil belajar dari data gabungan.} Alur
dalam federasi:
\begin{enumerate}
  \item (\emph{offline}, sekali) peneliti menerbitkan kamus SFM dari Paper~1 ke ranah publik;
  \item (tiap klien, lokal dan mandiri) menerapkan kamus ke datanya sendiri untuk
  memperoleh sembilan fitur SFM --- klien-CIC memakai pemetaan CICFlowMeter,
  klien-UNSW memakai pemetaan Argus, tanpa bertukar data;
  \item (federasi) kedua klien kini berada pada ruang sembilan dimensi yang
  \emph{sama secara semantik}, sehingga agregasi bobot menjadi sah.
\end{enumerate}
Yang \emph{tidak} dilakukan: melakukan \emph{fitting} pemetaan atau normalisasi
pada data gabungan kedua klien, maupun mengintip data mentah pihak lain.

\subsection{SFM sebagai \emph{interlingua} fitur dan \emph{adapter} per-alat}
\label{sec:adapter}

Secara lebih umum, yang dipetakan adalah skema \emph{alat} ekstraktor, bukan
datanya. SFM berperan sebagai \emph{interlingua}: tiap alat (CICFlowMeter,
Argus/Bro, NFStream, Zeek, \dots) diterjemahkan ke sembilan fitur SFM yang sama
melalui sebuah \emph{adapter} publik berisi nama kolom sumber, rumus derivasi, dan
catatan validasi. Klien baru dengan alat $Z$ cukup mengambil \emph{adapter}~$Z$
dari \emph{registry} publik, menerapkannya secara lokal, lalu berfederasi hanya
dalam ``bahasa'' sembilan fitur SFM. Server tidak perlu mengetahui alat klien
maupun melihat datanya.

Keterbatasan diakui secara jujur: bila alat klien belum memiliki \emph{adapter},
seorang pakar domain perlu membuatnya sekali; jadi SFM tidak sepenuhnya otomatis.
Lebih jauh, \emph{adapter} hanya sah bila alat sasaran benar-benar mengukur besaran
yang bersesuaian secara semantik --- bila sebuah besaran tak terukur (mis.\ byte
per-arah), fitur itu tak dapat dipetakan. Ini batas kelayakan yang kami nyatakan di
Bagian~\ref{sec:limitations}.

\subsection{Keputusan arsitektur: MLP ringan pada ruang SFM}
\label{sec:arch}

XGBoost yang menjadi \emph{backbone} Paper~1/2 tidak natural untuk FL karena
agregasi bobot pohon tidak langsung. Karena itu model lokal pada Paper~4 adalah
\emph{multilayer perceptron} (MLP) ringan di atas sembilan fitur SFM, dengan
arsitektur masukan~$9 \rightarrow [64, 32] \rightarrow 1$ (ReLU pada lapis
tersembunyi, \emph{sigmoid} pada keluaran). MLP ini memungkinkan FedAvg langsung
atas bobot. Sebagai \emph{jembatan} ke Paper~1, kami tetap melaporkan XGBoost
terpusat dan memastikan MLP terpusat tidak kalah jauh pada sembilan fitur SFM
(uji kewarasan). Arsitektur sengaja kecil: tujuan Paper~4 bukan mengejar MLP
terbaik, melainkan menguji SFM sebagai \emph{enabler} FL.

\subsection{Protokol federasi}
\label{sec:protocol}

Agregasi mengikuti FedAvg \eqref{eq:fedavg} dengan pembanding FedProx
\eqref{eq:fedprox} yang menambahkan suku \emph{proximal} untuk meredam
penyimpangan klien di bawah non-IID. Normalisasi \emph{z-score} dihitung
\emph{per-klien lokal} (opsi paling menjaga privasi); statistik tidak dibagikan.
Pertukaran parameter dilakukan melalui penyimpanan objek (S3), sehingga antar-node
tidak perlu membuka \emph{socket} langsung.

\begin{align}
  w_{t+1} &= \sum_{k} \frac{n_k}{n}\, w_k^{(t)} \label{eq:fedavg}\\[2pt]
  \min_{w}\; & F_k(w) + \frac{\mu}{2}\lVert w - w_{\text{global}} \rVert^2 \label{eq:fedprox}
\end{align}

\noindent di mana $n_k$ adalah ukuran data klien~$k$, $n=\sum_k n_k$, $F_k$ rugi
lokal (\emph{binary cross-entropy}), dan $\mu$ bobot \emph{proximal} FedProx.
Metrik utama adalah \emph{Matthews Correlation Coefficient} (MCC),
\begin{equation}
  \mathrm{MCC} = \frac{TP\cdot TN - FP\cdot FN}
  {\sqrt{(TP{+}FP)(TP{+}FN)(TN{+}FP)(TN{+}FN)}},
  \label{eq:mcc}
\end{equation}
dilengkapi \emph{recall}, \emph{precision}, F1, dan \emph{balanced accuracy}
(penting karena data tak-seimbang).

\begin{algorithm}[t]
\caption{FedAvg/FedProx di ruang SFM (per \emph{run})}
\label{alg:fed}
\begin{algorithmic}[1]
\State Server inisialisasi bobot global $w_0$; sebar ke semua klien.
\For{ronde $t=0,\dots,R-1$}
  \For{tiap klien $k$ (paralel)}
    \State Terapkan kamus SFM ke data lokal; \emph{z-score} lokal.
    \State $w_k^{(t)} \gets$ latih $E$ \emph{epoch} dari $w_t$ (FedProx bila $\mu>0$).
    \State Kirim $w_k^{(t)}$ dan ukuran $n_k$ (hanya bobot, bukan data).
  \EndFor
  \State $w_{t+1} \gets \sum_k (n_k/n)\, w_k^{(t)}$ \Comment{FedAvg}
\EndFor
\State \Return model global $w_R$.
\end{algorithmic}
\end{algorithm}

\subsection{Pengaturan eksperimen}
\label{sec:setup}

\emph{Dataset:} CSE-CIC-IDS2018 \cite{cicids} dan UNSW-NB15 \cite{unsw}, memakai
sembilan fitur SFM yang telah tervalidasi pada Paper~1 (tidak divalidasi ulang).
\emph{Skenario:} (A) dua klien $\{\text{CIC},\text{UNSW}\}$ dengan non-IID natural
akibat beda skema asli yang disatukan via SFM; (B) perluasan ke $K$ klien
(mis.\ $4$--$10$) melalui partisi non-IID berbasis label (Dirichlet). \emph{Baseline:}
terpusat (\emph{upper bound}), lokal-saja (\emph{lower bound}), dan federated-SFM
(FedAvg + FedProx), serta ablasi FL \emph{dengan} vs \emph{tanpa} SFM. \emph{Protokol:}
lima \emph{seed} $\{13,42,101,202,303\}$, rerata $\pm$ simpangan baku, selang
kepercayaan 95\%, dan uji signifikansi (pola Paper~2). Metrik komunikasi tambahan:
jumlah ronde hingga konvergen dan \emph{byte} per ronde terhadap MCC.

\paragraph{Infrastruktur.} Federasi dijalankan pada tiga instans EC2 ber-IP publik
(satu agregator, dua klien) di dalam satu VPC dengan \emph{Internet Gateway} tanpa
NAT, \emph{Security Group} ketat (tanpa SSH terbuka; akses operator via SSM), dan
pertukaran bobot melalui S3. Data mentah tiap klien tetap hanya berada di
node-nya. Logika agregator/klien yang dieksekusi di EC2 identik dengan yang
disimulasikan pada \emph{notebook} evaluasi, sehingga hasil simulasi memvalidasi
skrip sebelum penggelaran.

% ============================================================
\FloatBarrier
\section{Hasil}
\label{sec:results}

\begin{quote}
\small\itshape Catatan kejujuran: seluruh sel bertanda \TBD{} akan diisi dari
eksperimen nyata (notebook \texttt{nb01--nb07} + eksekusi AWS). Tidak ada angka
ilustratif pada draf ini.
\end{quote}

\subsection{Jembatan \emph{backbone} dan batas terpusat/lokal (H2)}
Tabel~\ref{tab:main} membandingkan terpusat, lokal-saja, dan federated-SFM pada
tiga himpunan uji. Urutan yang diharapkan: terpusat $\ge$ federated-SFM $\ge$
lokal-saja; besar selisih menunjukkan seberapa jauh FL menutup celah tanpa
memusatkan data.

\begin{table}[t]
\centering
\caption{MCC (rerata $\pm$ sb, lima \emph{seed}) pada CIC-\emph{test},
UNSW-\emph{test}, dan GLOBAL-\emph{test}. Nilai diisi dari \texttt{nb03}/\texttt{nb07}.}
\label{tab:main}
\begin{tabular}{lccc}
\toprule
Pengaturan & CIC-\emph{test} & UNSW-\emph{test} & GLOBAL-\emph{test} \\
\midrule
Terpusat (\emph{upper bound}) & \TBD & \TBD & \TBD \\
Lokal-saja (\emph{lower bound}) & \TBD & \TBD & \TBD \\
Federated-SFM (FedAvg) & \TBD & \TBD & \TBD \\
\bottomrule
\end{tabular}
\end{table}

\subsection{Ablasi: SFM sebagai \emph{enabler} (H1)}
Tabel~\ref{tab:ablation} menguji klaim inti: FedAvg di ruang SFM selaras versus
upaya naif tanpa SFM (penumpukan \emph{by-position}; irisan nama kolom). Bukti
dapat bersifat kuantitatif (MCC SFM jauh di atas naif) maupun kualitatif (bila
irisan nama kosong, federasi lintas-skema bahkan tak terbentuk tanpa SFM).

\begin{table}[t]
\centering
\caption{Ablasi FL dengan vs tanpa SFM, MCC GLOBAL-\emph{test}. Dari \texttt{nb04}.}
\label{tab:ablation}
\begin{tabular}{lcc}
\toprule
Varian & Dimensi fitur & MCC GLOBAL-\emph{test} \\
\midrule
Federated-SFM (9 selaras) & 9 & \TBD \\
Naif-1 (\emph{by-position}) & 9 & \TBD \\
Naif-2 (irisan nama) & \TBD & \TBD \\
\bottomrule
\end{tabular}
\end{table}

\subsection{Pengaruh non-IID dan jumlah klien (H3)}
Tabel~\ref{tab:noniid} melaporkan MCC global terhadap jumlah klien $K$ dan derajat
non-IID (parameter Dirichlet $\alpha$). Diharapkan: $\alpha$ kecil (non-IID
ekstrem) dan $K$ besar menurunkan MCC; besarnya penurunan dilaporkan apa adanya.

\begin{table}[t]
\centering
\caption{MCC GLOBAL-\emph{test} vs $(K,\alpha)$. Dari \texttt{nb05}.}
\label{tab:noniid}
\begin{tabular}{lccc}
\toprule
$K$ & $\alpha=0{,}1$ & $\alpha=1{,}0$ & $\alpha=100$ \\
\midrule
2  & \TBD & \TBD & \TBD \\
4  & \TBD & \TBD & \TBD \\
6  & \TBD & \TBD & \TBD \\
10 & \TBD & \TBD & \TBD \\
\bottomrule
\end{tabular}
\end{table}

\subsection{FedAvg vs FedProx dan biaya komunikasi}
Tabel~\ref{tab:comm} membandingkan FedAvg dan FedProx di bawah non-IID beserta
biaya komunikasi (jumlah parameter, \emph{byte} per ronde, ronde hingga konvergen).

\begin{table}[t]
\centering
\caption{FedAvg vs FedProx + biaya komunikasi. Dari \texttt{nb06}.}
\label{tab:comm}
\begin{tabular}{lcccc}
\toprule
Metode & MCC & Ronde-ke-konvergen & \emph{Byte}/ronde & Total \emph{byte} \\
\midrule
FedAvg          & \TBD & \TBD & \TBD & \TBD \\
FedProx ($\mu{=}0{,}001$) & \TBD & \TBD & \TBD & \TBD \\
FedProx ($\mu{=}0{,}01$)  & \TBD & \TBD & \TBD & \TBD \\
FedProx ($\mu{=}0{,}1$)   & \TBD & \TBD & \TBD & \TBD \\
\bottomrule
\end{tabular}
\end{table}

\subsection{Signifikansi statistik}
Uji berpasangan (\emph{t-test} dan Wilcoxon) atas lima \emph{seed} membandingkan
federated-SFM dengan lokal-saja; nilai $p$ dan ukuran efek (selisih rerata) diisi
dari \texttt{nb07}. Dengan lima \emph{seed}, daya uji Wilcoxon rendah; kedua uji
dilaporkan bersama ukuran efek.

% ============================================================
\FloatBarrier
\section{Diskusi}
\label{sec:discussion}

Jika hipotesis terkonfirmasi, temuan utama bukanlah ``FL bekerja untuk NIDS'' ---
itu sudah diketahui --- melainkan bahwa \emph{interoperabilitas fitur} adalah
prasyarat yang selama ini terlewat, dan SFM menyediakannya sebagai lapisan
\emph{public prior} yang murah dan menjaga privasi. Hasil ablasi (H1) menjadi
pembeda terkuat: tanpa jembatan semantik, penggabungan fitur lintas-\emph{tool}
tidak bermakna, bahkan dapat gagal terbentuk. Perbandingan terhadap batas terpusat
(H2) mengukur ``harga privasi'', sedangkan studi non-IID (H3) memetakan batas
praktis penerapan.

% ============================================================
\FloatBarrier
\section{Keterbatasan}
\label{sec:limitations}

Pertama, SFM membutuhkan \emph{adapter} per-alat; alat tanpa \emph{adapter}
memerlukan upaya pakar sekali, sehingga SFM tidak sepenuhnya otomatis. Kedua,
sebagian besaran mungkin tak terpetakan antar-alat (batas kelayakan semantik),
seperti dicontohkan pada Paper~1. Ketiga, eksperimen inti memakai dua
alat/dua klien sebagai instansiasi konkret; generalisasi ke lebih banyak alat
adalah kerja lanjutan. Keempat, kajian ini belum mencakup ketahanan adversarial
dan serangan \emph{poisoning} dalam setting federasi, yang dicadangkan untuk
makalah berikutnya.

% ============================================================
\FloatBarrier
\section{Kesimpulan dan Kerja Mendatang}
\label{sec:conclusion}

Kami mengajukan SFM sebagai \emph{enabler} \emph{federated learning} NIDS
lintas-skema: kamus pemetaan semantik publik memungkinkan klien berskema berbeda
berfederasi tanpa berbagi data mentah. Draf ini memaparkan motivasi, kerangka
privasi, arsitektur, protokol, dan rancangan evaluasi yang menyeluruh; seluruh
angka hasil akan diisi dari eksperimen nyata sebelum publikasi. Kerja mendatang
meliputi perluasan \emph{registry adapter}, integrasi ketahanan adversarial dan
pertahanan \emph{poisoning}, serta mekanisme normalisasi lintas-klien yang lebih
selaras namun tetap menjaga privasi.

% ============================================================
\section*{Ketersediaan Kode dan Data}
Implementasi node federasi (\texttt{aws/fed\_node.py}), templat infrastruktur
(\texttt{aws/fed-vpc-3ec2-public.yaml}), \emph{runbook}, dan \emph{notebook}
eksperimen (\texttt{nb01--nb07}) disertakan pada repositori proyek. Kamus SFM
diwarisi dari Paper~1. Dataset sumber bersifat publik \cite{cicids,unsw}.

\section*{Pernyataan Etika}
Seluruh trafik serangan yang dipakai berasal dari \emph{dataset} publik; eksekusi
federasi dilakukan di lingkungan \emph{cloud} terisolasi milik sendiri. Tidak ada
data pribadi yang dikumpulkan.

% ============================================================
\begin{thebibliography}{99}
% Format IEEE. Nomor urut sesuai kemunculan; label dari federated_refs_clean.md.

\bibitem{sfm} H.~Y.~Martono, I.~Syarif, and F.~A.~Saputra, ``Cross-dataset
generalization of ML-based network intrusion detection via semantic feature
mapping and few-shot adaptation, validated on real cloud traffic,'' \emph{status
terbit diisi saat accepted}.

\bibitem{popoola} S.~I.~Popoola, G.~Gui, B.~Adebisi, M.~Hammoudeh, and
H.~Gacanin, ``Federated deep learning for collaborative intrusion detection in
heterogeneous networks,'' in \emph{Proc. IEEE 94th Veh. Technol. Conf.
(VTC-Fall)}, 2021, pp.~1--6.

\bibitem{fsurvey1} A.~Khraisat, A.~Alazab, T.~Jan, and A.~Gopez, ``Survey on
federated learning for intrusion detection system: Concept, architectures,
aggregation strategies, challenges, and future directions,'' \emph{ACM Comput.
Surv.}, 2024.

\bibitem{fsurvey2} J.~L.~Hern\'andez-Ramos \emph{et~al.}, ``Intrusion detection
based on federated learning: A systematic review,'' \emph{ACM Comput. Surv.},
2025.

\bibitem{ftaxonomy} A.~A.~Wardana and P.~Sukarno, ``Taxonomy and survey of
collaborative intrusion detection system using federated learning,'' \emph{ACM
Comput. Surv.}, 2024.

\bibitem{falgo} E.~V.~Fedorchenko, E.~Novikova, and A.~N.~Shulepov, ``Comparative
review of the IDS based on federated learning,'' \emph{Algorithms}, vol.~15,
no.~7, art.~247, 2022.

\bibitem{fresource} R.~Doriguzzi-Corin, S.~Cretti, and D.~Siracusa,
``Resource-efficient federated learning for network intrusion detection,'' in
\emph{Proc. IEEE NetSoft}, 2024.

\bibitem{fdbscan} Y.-C.~Lee, W.-C.~Chien, and Y.-C.~Chang, ``FedDB: A federated
learning approach using DBSCAN for DDoS attack detection,'' \emph{Appl. Sci.},
2024.

\bibitem{fentente} J.~Xu, C.~Li, Y.~Zheng, and Z.~Li, ``Entente: Cross-silo
intrusion detection on network log graphs with federated learning,''
\emph{arXiv:2503.14284}, 2025.

\bibitem{fiiot} J.~Mao, Z.~Wei, B.~Li, R.~Zhang, and L.~Song, ``FedIn-NID: A FL
framework for NID in large-scale heterogeneous industrial IoT,'' \emph{IEEE Trans.
Inf. Forensics Security}, 2025.

\bibitem{fnoniid} L.~Lavaur, Y.~Busnel, and F.~Autrel, ``Highlighting the limits
of federated learning in intrusion detection,'' in \emph{Proc. IEEE ICDCS}, 2024.

\bibitem{fpersonal} S.~Gurpreet, S.~Keshav, P.~Rajalakshmi, and Y.~Xong,
``Sentinel: Dynamic knowledge distillation for personalized federated intrusion
detection in heterogeneous IoT networks,'' \emph{arXiv:2510.23019}, 2025.

\bibitem{fmist} S.~Ali, R.~Alireza, and A.~Mahmood, ``Mist-assisted federated
learning for intrusion detection in heterogeneous IoT networks,''
\emph{arXiv:2511.00271}, 2025.

\bibitem{flimit} L.~Lavaur, Y.~Busnel, and F.~Autrel, ``Demo: Highlighting the
limits of federated learning in intrusion detection,'' in \emph{Proc. IEEE
ICDCS}, 2024.

\bibitem{fcrossvan} M.~Khalid, U.~Zukaib, M.~Al-Rakhami, and A.~M.~Alamri,
``FedCross-VAN: Federated cross-domain behavior alignment for VANET intrusion
detection,'' \emph{Trans. Emerg. Telecommun. Technol.}, vol.~37, 2025.

\bibitem{frobust} J.~Zhang \emph{et~al.}, ``Delving into the adversarial
robustness of federated learning,'' \emph{arXiv:2302.09479}, 2023.

\bibitem{fbackdoor} Z.~Sun, P.~Kairouz, A.~T.~Suresh, and H.~B.~McMahan, ``Can you
really backdoor federated learning?'' \emph{arXiv:1911.07963}, 2019.

\bibitem{fdelve} V.~Rey, P.~M.~S.~S\'anchez, A.~H.~Celdr\'an, G.~Bovet, and
M.~Jaggi, ``Federated learning for malware detection in IoT devices,''
\emph{Comput. Netw.}, 2021.

\bibitem{ffewshot} Y.~Hu, J.~Wu, G.~Li, J.~Li, and J.~Cheng, ``Privacy-preserving
few-shot traffic detection against APTs via federated meta learning,'' \emph{IEEE
Trans. Netw. Sci. Eng.}, vol.~11, pp.~2549--2560, 2024.

\bibitem{ffecograph} Q.~Mao, X.~Lin, G.~Li, and J.~Li, ``FeCoGraph: Label-aware
federated graph contrastive learning for few-shot network intrusion detection,''
\emph{IEEE Trans. Inf. Forensics Security}, 2025.

\bibitem{ftransfer} J.~Zhang, C.~Luo, M.~Carpenter, and G.~Min, ``Federated
learning for distributed IIoT intrusion detection using transfer approaches,''
\emph{IEEE Trans. Ind. Informat.}, vol.~19, pp.~8159--8169, 2023.

\bibitem{sarhan} M.~Sarhan, S.~Layeghy, and M.~Portmann, ``Towards a standard
feature set for network intrusion detection system datasets,'' \emph{Mobile Netw.
Appl.}, vol.~27, no.~1, pp.~357--370, 2022.

\bibitem{sarhannf} M.~Sarhan, S.~Layeghy, N.~Moustafa, and M.~Portmann, ``NetFlow
datasets for machine learning-based network intrusion detection systems,'' in
\emph{Big Data Technologies and Applications}. Cham, Switzerland: Springer, 2021.

\bibitem{cicids} I.~Sharafaldin, A.~H.~Lashkari, and A.~A.~Ghorbani, ``Toward
generating a new intrusion detection dataset and intrusion traffic
characterization,'' in \emph{Proc. 4th ICISSP}, 2018, pp.~108--116.

\bibitem{unsw} N.~Moustafa and J.~Slay, ``UNSW-NB15: A comprehensive data set for
network intrusion detection systems,'' in \emph{Proc. MilCIS}, 2015, pp.~1--6.

\end{thebibliography}

\end{document}
